\par\SourceLeaf{140}
Da dieses Gutachten und die ihm gegebene Auslegung die eigentliche Grundlage der Einrichtung bilden, die den Hauptgegenstand der 1. Abteilung meiner Denkschrift ausmacht, nämlich der Progressionsnorm, werde ich es weiter unten einer ausführlichen und genauen Prüfung unterziehen und beschränke mich deshalb hier vorerst auf folgende Bemerkungen.

Aus den Akten geht nicht hervor, weshalb \textsc{Kritter} am Anfang seines Gutachtens von 22—24 verheirateten Professoren spricht, da P. in seinem Schreiben ausdrücklich 26 und nur diese Zahl genannt hatte. Ich vermute aber, dass \textsc{Kritter} jene Zahlen 22—24 als für eine \textit{frühere} Zeit gültig angenommen hat. Meines Wissens sind jedoch keine vollständigen Register über die persönlichen Verhältnisse der Mitglieder der Witwenkasse so geführt worden, dass sich daraus für jeden beliebigen Zeitpunkt die Anzahl der \textit{Verheirateten} entnehmen ließe. Entweder hat \textsc{Kritter} jene Zahlen also nur aus der Zahl aller Teilnehmer in früherer Zeit durch eine ungefähre Schätzung \textit{erschlossen}, oder sie beruhen auf besonderen Mitteilungen, die sich dann der Natur der Sache nach nur auf Zeitpunkte beziehen können, die nicht viele Jahre zurücklagen und noch im Gedächtnis lebendig waren.

\textsc{Kritter}s Antwort auf die erste Frage besteht dann kurz darin, dass man für die \textit{Zeit}, zu der das Maximum eingetreten sei, eine Witwe auf etwa zwei bestehende Ehen rechnen könne, also für jene 22—24 Ehen 12 Witwen. Dies habe sich, wie er angibt, nach dem Durchschnitt der letzten 17 Jahre insofern bestätigt, als es bald mehr, bald weniger als 12 Witwen gegeben habe. Dass dann die der damaligen Zahl von 26 Ehen entsprechende Witwenzahl um eine größer sein würde, ist nicht ausdrücklich gesagt, aber darin enthalten. Auf die zweite Frage gibt er an, dass man nach den Erfahrungen der Bremer Witwenkasse, wo die Waisenpensionierung nur bis zum \(18^{\text{ten}}\) Lebensjahr dauere, mit einer Vergrößerung der Pensionenzahl um den sechsten Teil rechne; bei der hiesigen Kasse, wo die Dauer 2 Jahre länger sein solle, müsse man also wohl etwas mehr annehmen. — Bei den 24 Ehen kommen wir demnach auf etwas mehr als 14, bei den 26 Ehen, nachdem ihre Wirkung vollständig eingetreten ist, auf 15 nach \textsc{Kritter}s Worten oder auf etwas mehr als \(15\tfrac12\), das heißt auf beinahe 16, nach den von ihm ausgesprochenen Grundsätzen.

Ob man, \textit{ganz abgesehen von der näheren Prüfung des Inhalts des Gutachtens}, eine solche Behandlung des Gegenstands eine \textit{ausführliche} und \textit{genaue} Untersuchung nennen könne, wie \textsc{Kaestner} sie für notwendig gehalten hatte, lasse ich hier auf sich beruhen. P. entwarf nun aber auf der Grundlage dieses Gutachtens ein P. M., worin er zeigt, dass die vom Kuratorium vorgeschlagene Erhöhung der Pensionen, wenn sie ohne weitere besondere Erhöhung für die sechs ältesten Witwen für alle einheitlich auf 130 Rth. festgesetzt werde, ohne jede Gefahr für die Kasse auch dann erfolgen könne, wenn die jährlichen Beiträge nicht erhöht würden. Werde aber die Erhöhung der Beiträge in der vom Kuratorium angegebenen Weise angenommen, so könne auch die Verlängerung der Waisenpensionen umso sicherer eingeführt werden, weil unter den gegebenen Umständen nicht zu erwarten sei, dass diese Änderung bald wirksam werde. P. schließt schließlich seinen Vortrag, dem ich \textit{bis hierher} ohne Bedenken meinen vollen Beifall gebe, mit folgendem kurzen Zusatz. Da hier \textit{zum ersten Mal} der Gegenstand meiner eigenen Untersuchung, nämlich die Progressionsnorm, auf den Schauplatz tritt, schreibe ich ihn vollständig und getreu mit P.s eigenen Worten ab:

\begin{quote}
‘Bei der allgemeinen Erhöhung der Pensionen auf 130 Rth. scheint mir nicht die geringste Gefahr zu bestehen:
\end{quote}
\SourceLeaf{141}
\begin{quote}
Vielmehr scheint es mir noch
\begin{list}{}{\setlength{\leftmargin}{1.8em}\setlength{\labelwidth}{1.4em}\setlength{\labelsep}{.4em}}
\item[3)] möglich und dienlich, dass es jetzt zur dauernden Norm gemacht werden dürfe, die Witwenpensionen jedes Mal um 10 Thl. zu erhöhen, sobald sich das Grundvermögen um 5000 Thl. vermehrt hat und die Zahl der Witwen noch nicht über das Maximum von 15 gestiegen ist. Es ist klar, dass man dies tun kann, denn eine Erhöhung von 10 Thl. für 15 Witwen beträgt 150 Thl., und 5000 Thl. zu 3 Prozent bringen ebenso viel Zinsen. Dass aber der nächste Erhöhungstermin bald eintreten kann, auch wenn unsere Kasse keine außerordentlichen Zuflüsse erhält, lässt sich wenigstens als sehr wahrscheinlich berechnen. Da wir gegenwärtig nur 10 Witwen zu pensionieren haben, müssen, wenn der neue Zuschlag zu den Beiträgen bewilligt wird, jedes Jahr über 1000 Thl. in der Kasse bleiben; folglich müssen 5000 Thl. schon in 5 Jahren zum Grundvermögen hinzugekommen sein, wenn sich die Zahl der Witwen inzwischen nicht vermehrt. Nimmt man aber auch den höchst unwahrscheinlichen Fall an, dass die Zahl jedes Jahr um eine Witwe zunähme, bis sie das Maximum von 15 erreicht hätte, so könnte es doch kaum 10 Jahre dauern.
\end{list}
Gut wäre es wenigstens, wenn dieser Umstand im Bericht an die Königliche Regierung erwähnt würde.’
\end{quote}

\noindent\hspace*{\parindent}Da das Nötige über den wesentlichen Inhalt dieses Artikels weiter unten bei der Prüfung des \textsc{Kritter}schen Gutachtens und der daran geknüpften Folgerungen sowie an anderen Stellen vorkommen wird, sollen hier nur ein paar Nebenumstände berührt werden.

I. Auffallend ist aus der Feder des Kurators der Witwenkasse die \textit{unrichtige} Angabe der Witwen- oder vielmehr Pensionenzahl. Damals gab es nicht 10 Witwen und auch nicht 10 Pensionen, sondern 8 Witwen und, unter Hinzurechnung der Kinder des am 8. Juni 1794 verstorbenen Professors B., zusammen 9 Pensionen. Diese Unrichtigkeit lässt sich auch nicht etwa damit erklären, dass der Kurator damals noch nichts vom Wegfall der zuletzt auswärts (in Halle) verstorbenen Witwe (M.) gewusst habe. Denn es findet sich, dass der Betrag der Pension für das letzte halbe Jahr (Michaelis 1793 bis Ostern 1794) gegen eine von P. mitunterzeichnete und auf den 3. Mai 1794 datierte Quittung der Erbin ausgezahlt worden ist.

II. Die Schlusszeile (‘Gut wäre es wenigstens usw.’) lässt uns hinsichtlich der Frage etwas im Dunkeln, wofür der abgeschriebene Artikel eigentlich gelten soll: für einen förmlichen, zur Beschlussfassung gestellten Antrag oder nur für eine hingeworfene Idee. Weder im Reskript noch in den vorhergehenden Verhandlungen war von einer bestimmten, für alle künftigen Zeiten geltenden Normierung der veränderlichen Pensionshöhe die Rede gewesen. Es war dies also ein \textit{vierter}, zu den bereits beratenen neu hinzukommender Gegenstand, und zwar einer, der den drei anderen an Wichtigkeit keineswegs nachstand und die sorgfältigste Prüfung von allen Seiten erforderte. Wer einen eigens auf ein solches Ziel gerichteten Antrag einbringt, ist sich doch der Wichtigkeit der Sache bewusst; diese tritt aber durch die Bezeichnung einer solchen Lebensfrage als ‘dieser Umstand’ schwerlich genügend hervor. Bei einer Äußerung hingegen, die nur den Charakter eines gelegentlich hingeworfenen Gedankens hat, ist man schon nachsichtiger gegenüber einer durch Unachtsamkeit unterlaufenen Unrichtigkeit und gegenüber einer noch mangelhaften Formulierung. Man lasse sich hierbei durch die Nummerierung mit (3) nicht irremachen. P. \textit{zählt} in seinem Aufsatz nicht die Vorschläge, sondern die Folgerungen, die er aus den vorausgeschickten überschlägigen Rechnungen ableitet. Seine Nr. 1 enthält \textit{die} Folgerung, dass sich alle gegen die Verlängerung der Waisenpensionen vorgebrachten oder noch vorzubringenden Einwände erledigen, wenn man die vorgeschlagene Erhöhung der Beiträge annimmt; und seine Nr. 2 die Folgerung, dass die allgemeine Erhöhung der Pensionen auf 130 Rth. zwar ohne Weiteres sofort erfolgen könne, die besondere Erhöhung auf 150 Rth. für die sechs ältesten Witwen hingegen weder gerecht noch ratsam sei.

